%% ================================================================
%%  8. DISCUSSION: COORDINATE LOCALIZATION INSIGHTS
%%
%%  Cross-reference to the H-*/O-* taxonomy
%%  (Section~\ref{sec:taxonomy}, operational modes at
%%  Section~\ref{sec:operational-failures}) sits at the top of the section.
%% ================================================================
\section{Discussion: Coordinate Localization Insights}
\label{sec:discussion}

This section discusses three coordinate-localization findings
arising from the v6/v7 benchmark campaigns: the negative result on
text-label snapping (Stage~B), the deterministic-routing
contribution of Stage~C Phase~A, and the variance-vs-peak-precision
tradeoff between $N{=}1$ and $N{=}3$ Phase~1 aggregation.  For the
cognitive-versus-operational failure-mode taxonomy that grounds
some of the references below, see Section~\ref{sec:taxonomy} and
the operational modes at Section~\ref{sec:operational-failures}.

\subsection{Stage B: A Negative Result on Text-Label Snapping}
\label{sec:discussion-stage-b}

Between the v6 and v7 campaigns we added a \textbf{Stage~B}
refinement layer: after Phase~1 produces a coarse coordinate
estimate, we crop a $120 \times 120$\,px window around the
estimate, run Windows' built-in OCR engine, and---when the target
description names a specific text label---snap the estimate to the
center of the matching OCR bounding box.  The hit condition is
deliberately conservative: the Phase~1 coordinate must lie
\emph{inside} the matched bounding box.  This condition guarantees
that a successful snap never moves the estimate further from
ground truth than Phase~1 alone, but it also narrows the
applicability of the method to UI elements whose click-center is
inside the text-label bbox.

The v7 measurement outcome is \textbf{0 hits in 50 cases}
(Table~\ref{tab:stage-b-scope}).  The miss population partitions
cleanly by category.  Three categories (\texttt{ui\_dense},
\texttt{natural}, \texttt{precision}) contain no quoted text label
in the target description, so the label-extraction preprocessor
short-circuits to a miss before the OCR probe even runs.  The
remaining two categories (\texttt{ui\_button}, \texttt{ambiguous})
do contain quoted labels, but the OCR-detected text bounding
boxes are systematically offset from the intended click-centers by
$\sim$22--29\,px on real toolbar buttons: the ground-truth click
target is the button's geometric center, while the text label
occupies a left- or padding-offset region inside the button.
Phase~1 therefore does not land inside the label bbox, and the
conservative hit condition correctly refuses to snap.

We report this negative result as a \emph{scope-definition}
finding rather than as a failure.  The contract ``a snap never
makes the result worse'' was preserved: miss-path Pinpoint is
byte-identical to the v6 pipeline, and the overall v7 metrics
still improve substantially over v6 (mean $54.3 \to 29.9$\,px,
$\leq$\,5\,px $8.2\% \to 22.0\%$; Table~\ref{tab:results-guardrail-mv}).
Those improvements come from Phase~1 and Phase~2 stability on a
second independent run, not from Stage~B.

The result also tightens the specification for a successor layer.
A post-Phase~1 refinement layer that aspires to beat Phase~2 on
button targets cannot rely on OCR-text centers; it must either
recover the \emph{button geometry} (so that the snap target is the
enclosing rectangle's center, not the label's) or use a
geometrically aware feature extractor such as an edge / contour
detector.  This successor, \textbf{Stage~C (Semantic--Geometric
Refinement)}, is implemented and evaluated in the next subsection
(Section~\ref{sec:discussion-stage-c}).  In that framing, Stage~B
becomes the \texttt{menu\_item} sub-case of Stage~C---the one case
where text-label snapping is geometrically aligned with the click
target.

\subsection{Stage~C Phase~A: A Deterministic Routing Layer That
  Exposes Upstream Variance}
\label{sec:discussion-stage-c}

Stage~C replaces the single bbox-containment test of Stage~B with a
dispatcher over \texttt{target\_type}-specific CV pipelines.
Phase~1.5 emits a structured declaration
$\{\texttt{target\_type}, \texttt{intent}, \texttt{label\_hint}\}$
via a rule-based heuristic over the target prompt, and Stage~C
dispatches to the matching pipeline: \texttt{button} or \texttt{icon}
(Canny + morphological close +
\texttt{findContours} + aspect-ratio/area filter + nearest-centroid
selection, with a piecewise-linear distance-to-confidence map),
\texttt{menu\_item} (the Stage~B OCR-snap helper, retrofitted), or
\texttt{generic} (a miss-only passthrough that returns Phase~1
coordinates unchanged, preserving measurement honesty in the paper
ablation).  Only the \texttt{menu\_item} and contour pipelines are
VLM-touching in any form; \texttt{generic} never invokes OCR or CV
operators.  A rule-based \texttt{declare\_target} is used instead
of a VLM tool-use call so that routing variance is zero and any A/B
effect on hit rate is attributable to the CV pipeline, not to a
second VLM sampling step; a tool-use variant with the same contract
is left as follow-up work.  An implementation detail worth flagging
for future extension: the declaration is emitted in a nested
\texttt{phase\_1\_5\_structured} field in the benchmark result
schema; a later tool-use variant will want two of the embedded
fields (\texttt{source} and \texttt{confidence}) promoted to the
top level so a 3-arm comparison (rule-based / tool-use
deterministic / tool-use sampled) can group on them directly (the
three-arm protocol itself is described in the Phase~B roadmap at
the end of this subsection).

\paragraph{Stage~1 subset observation: a bimodal precision effect.}
A 20-case smoke (\texttt{ui\_button} + \texttt{ui\_dense},
\texttt{--repeat 3}) run as Stage~1 of the A/B protocol exhibited a
striking \emph{bimodal} split.  The left tail improved sharply
($\leq$\,5\,px $10\% \to 20\%$, $\leq$\,10\,px $20\% \to 33\%$,
$\leq$\,20\,px $50\% \to 58\%$) while the right tail degraded at the
same time ($\leq$\,50\,px $90\% \to 77\%$).  Inspecting the raw
per-repeat results makes the mechanism legible: Stage~C's button
pipeline delivers $\leq$\,1\,px hits on cases where Phase~1 lands
inside the button, but on the small number of cases where Phase~1
lands on a neighbouring button, Stage~C snaps to that neighbour's
centroid, moving the error from ``nearly right'' into
``deterministically wrong''.  The precision boost and the
catastrophic snap are two ends of the same deterministic rule.

\paragraph{Stage~2 full scale: the bimodal split does not survive
  single-sample evaluation.}
The apparent divergence between Stage~1 and Stage~2 is a sampling
artefact of the same deterministic rule: multi-sample averaging under
\texttt{--repeat 3} reveals both bimodal ends simultaneously, while
single-sample \texttt{--repeat 1} resolves each target onto one end
determined by Phase~1 routing.
At the 50-case, \texttt{--repeat 1} scale
(Table~\ref{tab:results-stage-progression}, v8 column), the Stage~1 bimodal pattern
disappears.  $\leq$\,5\,px drops $22.0\% \to 14.0\%$, $\leq$\,20\,px
drops $54.0\% \to 38.0\%$, overall $\leq$\,50\,px drops
$82.0\% \to 76.0\%$.  The per-category decomposition
(Table~\ref{tab:results-by-category}, v8 column) is what makes this result
interpretable rather than alarming: the $-60$\,pt collapse on
\texttt{precision} and the $-10$\,pt on \texttt{natural} come from
categories in which Stage~C never fires (10/10
\texttt{generic\_miss} each), so those deltas measure Phase~1
sampling variance between two independent runs, not Stage~C.
Conversely, the categories in which Stage~C \emph{does} route
through a specialized pipeline split by mechanism:
\texttt{ui\_button} ($+10$\,pt) is the path-A success case---Stage~C
reaches a button centroid when Phase~1 lands on the right
button---while \texttt{ui\_dense} ($-30$\,pt) is the path-B failure
case, where the icon pipeline's filter window is wide enough that
Phase~1 variance can flip the chosen contour from the intended icon
to a neighbour.  The Stage~1 bimodal observation is therefore real
but emergent only under multi-sample averaging; a single-sample
evaluation is dominated by whichever side of the amplifier Phase~1
happens to land on.

\paragraph{Cross-run reproducibility isolates the variance source.}
The strongest diagnostic is the Stage~1 $\leftrightarrow$ Stage~2
cross-run check on the $60$ rows where Stage~C dispatched to a
specialized pipeline.  Four wrong-detection hits---including
\texttt{ui\_button\_01} at exactly $(79, 22)$ for $95.0$\,px error
across all three Stage~1 repeats \emph{and} the Stage~2 single
run---are bit-identical.  Correspondingly, three perfect-hit cases
(\texttt{ui\_button\_02}, \texttt{\_04}, \texttt{\_08} at
$\leq$\,1\,px) are also bit-identical across runs.  Stage~C is
therefore a \emph{deterministic function of its inputs}; all
remaining across-run variance lives in the input
(\texttt{approx\_xy} from Phase~1).  Quantitatively, the case-level
standard deviations $\sigma_{\text{input}}$ and
$\sigma_{\text{output}}$ of the coordinates entering and leaving
Stage~C are essentially identical ($\sigma_{\text{in}} \approx
\sigma_{\text{out}}$ on both axes in Stage~1), confirming at the
distribution level what the bit-identical examples confirm at the
individual level: Stage~C neither attenuates nor amplifies variance
overall; it is a pure routing layer.  The unit-level CV-only
determinism suite (N=10 synthetic repeats, $\sigma < 0.5$\,px per
pipeline) corroborates this at the algorithmic layer: the empirical
bit-identical observations on real images reflect an underlying
algorithmic invariant verified at the unit-test level, so the
deterministic-function claim rests on three triangulating sources
(unit-level synthetic, case-level bit-identical, distribution-level
$\sigma$ preservation) rather than any single one.

\paragraph{What Stage~C's Phase~A actually establishes.}
The v8 numbers would, read naively, suggest that Stage~C ``made
things worse''---the overall $\leq$\,20\,px rate did fall by
$16$\,pt.  The per-category breakdown and the bit-identical
cross-run evidence refuse that reading.  Stage~C Phase~A
establishes three claims that together are stronger than any
left-tail success rate would have been: (i) a post-Phase~1
refinement layer can be run as a deterministic function, with no
routing variance of its own, with overall regression well below the
$+20$\,px threshold agreed \emph{a priori}; (ii) when that
refinement layer fires on the right input, it can reach the
geometric centroid of a button to within one pixel, cleanly
beating Phase~2 Monte Carlo voting on that subset; and (iii)
because Stage~C itself is deterministic, all remaining error,
including the Stage~1 $\leftrightarrow$ Stage~2 divergence and the
per-category asymmetry, localises unambiguously to Phase~1's VLM
sampling variance.  Future left-tail improvements therefore have
to address Phase~1, not Stage~C.

The three-agent organisation under which these claims were
produced---a supervisor, an implementer, and an independent
evaluator---is complementary to the task-phase decomposition
(planner / generator / evaluator) documented contemporaneously for
Anthropic's own harness~\cite{anthropic2026threeagentharness}.  Our
axis is \emph{responsibility-domain} based: the supervisor owns
design decisions and the decision log, the implementer owns the
code landing surface, and the evaluator owns the benchmark
artefacts and GO / NO-GO judgments.  The task-phase axis
decomposes a single reasoning episode; the responsibility-domain
axis decomposes the ownership of long-running artefacts.  Taken
together they suggest a two-dimensional design space for
multi-agent harnesses worth systematic exploration, with the
Phase~A evidence here contributing a datapoint in which the
responsibility-domain decomposition was sufficient to converge on
a paper-grade contribution inside a two-week implementation
window.

\paragraph{Phase~B roadmap.}
With the variance bottleneck localised, three follow-ups compose
cleanly.  First, the \texttt{ui\_dense} $-30$\,pt outcome pushes
the icon pipeline's filter window and search-radius tuning into
Phase~B scope; the dry-run evidence (Stage~1 per-case logs) points
at dense-toolbar geometries that the current aspect-ratio /
area cuts do not cover.  Second, and the clear priority from
Stage~2, is \emph{Phase~1 upstream variance reduction}: either
majority voting or bootstrap aggregation over $k \geq 3$
independent Phase~1 draws, or a seeded-sampling investigation at
the VLM API surface.  The cross-run reproducibility data suggest
that a single-sample Phase~1 spends its variance budget unevenly
across categories---\texttt{precision} drew the unlucky sample in
our v8 run---and that averaging a small number of draws should
preserve Stage~C's deterministic precision gain while filling in
its missing robustness.  Third, the rule-based \texttt{declare\_target}
admits a drop-in tool-use variant with the same schema; comparing
rule-based, deterministic-tool-use, and sampled-tool-use routing in
a three-arm smoke would quantify how much of the observed Stage~C
effect depends on the routing signal being correct (addressable by
a VLM) versus on its being stable (already achieved by the rule
base).  This three-arm design isolates three independently varying
factors---(i) routing-rule complexity (rule vs.\ tool-use),
(ii) VLM-variance exposure (deterministic vs.\ sampled), and
(iii) downstream Stage~C effect (held constant)---so any observed
change can be attributed to a single axis.  We position these as Phase~B; the Phase~A closeout is that
the deterministic routing layer itself is built and measured.

The v9 campaign (Section~\ref{sec:results}, ``v9 extension''
paragraph) implements the first two of these three follow-ups
(Workstream~1: $N=3$ majority voting at Phase~1; Workstream~2:
icon-pipeline adaptive tuning) and closes the prediction with a
measured $+6$\,pt gain on the primary $\leq$\,20\,px metric and a
$+70$\,pt recovery on \texttt{precision}, the category we flagged
above as the strongest Phase~1-VLM-variance evidence.  The third
follow-up---the rule~/~deterministic-tool-use~/~sampled-tool-use
three-arm smoke---remains Phase~C work.  The fact that Phase~B
delivered the Strong-GO outcome using only the two
upstream-variance + pipeline-conservativism follow-ups, without
yet varying the routing signal, corroborates the Phase~A
localisation claim: Phase~1 sampling variance was indeed the
dominant bottleneck, and reducing it was sufficient to produce the
predicted gain.


\subsection{N=1 vs N=3 Aggregation: A Variance-vs-Peak-Precision Tradeoff}
\label{sec:aggregation-discussion}

While arm~A and arm~B share identical model, prompts, and
Stage~C deterministic routing pipeline, they differ only in
the aggregation strategy applied at the phase-1 sample stage.
The resulting distributional contrast on the precision
category is informative not as evidence of model-specific
bias---Welch's $t$-test fails to reject equality of means at
$p=0.41$---but as evidence of aggregation strategy's
variance-reduction effect.

We conducted a paired benchmark comparison between two
aggregation strategies on the precision category ($n=10$
sub-pixel and tiny-target cases) using Claude Opus~4.7:
\textbf{arm~A} uses $N=3$ phase-1 sample median aggregation,
\textbf{arm~B} uses $N=1$ single-sample baseline.  Both arms
share identical model, prompts, target images, and Stage~C
deterministic routing pipeline; the only variation is the
aggregation strategy.

\begin{table}[t]
  \centering
  \small
  \caption{Per-case precision pinpoint distance (px) for
    arm~A ($N=3$ median) versus arm~B ($N=1$).
    \textbf{Bold} entries flag $>50$~px gross-error outliers
    in arm~B (absent from arm~A).}
  \label{tab:aggregation-percase}
  \resizebox{\columnwidth}{!}{%
  \begin{tabular}{@{}llrr@{}}
    \toprule
    \textbf{Case} & \textbf{Target} & \textbf{arm~A} & \textbf{arm~B} \\
    \midrule
    precision\_00 & single red pixel                         & 18.2 & \textbf{80.4} \\
    precision\_01 & thin blue needle (tip)                   & 1.0  & 4.0 \\
    precision\_02 & tiny green dot                           & 34.5 & \textbf{65.3} \\
    precision\_03 & small orange circle                      & 23.5 & 34.1 \\
    precision\_04 & crosshair intersection                   & 10.0 & 3.0 \\
    precision\_05 & purple diamond                           & 19.7 & \textbf{70.8} \\
    precision\_06 & top-left corner of white square          & 21.5 & 1.0 \\
    precision\_07 & right endpoint of horizontal yellow line & 20.2 & 2.2 \\
    precision\_08 & innermost ring of target                 & 6.0  & 12.8 \\
    precision\_09 & nose of small arrow                      & 28.4 & 1.4 \\
    \midrule
    \multicolumn{2}{r}{\textit{mean / stdev}}                & 18.30 / 10.15 & 27.50 / 32.55 \\
    \bottomrule
  \end{tabular}%
  }
\end{table}

\paragraph{Descriptive statistics and statistical tests.}
The mean of arm~B ($27.5$~px) is higher than arm~A ($18.3$~px),
but the medians invert: arm~B median is $8.4$~px while arm~A
median is $19.95$~px.  The standard deviations diverge sharply,
$32.55$~px in arm~B versus $10.15$~px in arm~A, a ratio of
$3.21\times$.  A Welch's two-sample $t$-test on the means
(unequal variances) yields $t=-0.853$, $df=10.73$, two-tailed
$p\approx 0.41$---i.e.\ the mean difference is \emph{not}
statistically significant, the variance inflation in arm~B
masking any central-tendency shift.  Conversely, an $F$-test
for variance equality yields $F=10.29$ on $(9,9)$ degrees of
freedom against $F_{\text{crit}}(\alpha=0.05)\approx 4.03$,
giving $p\approx 0.0019$---i.e.\ the variance difference is highly
significant.\footnote{Welch's $t$ failing to reject $H_0$ at
$p=0.41$, combined with $F$-test rejection at $p \approx 0.0019$,
formalizes the claim that $N=3$ median aggregation operates
on the variance dimension while leaving central tendency
invariant.}

\paragraph{Threshold-stratified hit rate.}
The hit-rate divergence by accuracy threshold makes the
distributional contrast practically interpretable:

\begin{center}
\small
\begin{tabular}{@{}lrrr@{}}
\toprule
\textbf{Threshold} & \textbf{arm~A hit} & \textbf{arm~B hit} & \textbf{$\Delta$} \\
\midrule
$\leq 5$~px  & 1/10 (10\%)  & \textbf{5/10 (50\%)} & $+40$\,pp arm~B \\
$\leq 10$~px & 3/10 (30\%)  & 5/10 (50\%)          & $+20$\,pp arm~B \\
$\leq 20$~px & 5/10 (50\%)  & 6/10 (60\%)          & $+10$\,pp arm~B \\
$\leq 50$~px & \textbf{10/10 (100\%)} & 7/10 (70\%) & $-30$\,pp arm~A \\
\bottomrule
\end{tabular}
\end{center}

\paragraph{Bimodal signature interpretation.}
The mean-versus-median divergence in arm~B (mean $27.5$,
median $8.4$)---combined with arm~A's mean-median agreement
($18.3$ vs.\ $19.95$)---is the textbook quantitative
signature of bimodal distribution.  Specifically, arm~B
exhibits a ``hit-or-miss'' pattern: $5/10$ cases achieve
sub-$5$~px precision (single-sample lucky), while $4/10$
cases produce $>34$~px gross errors (single-sample unlucky).
The $N=3$ median aggregation in arm~A trades this peak
precision (only $1/10$ $\leq 5$~px hits) for distributional
stability ($10/10$ $\leq 50$~px hits, zero gross-error
outliers).

\paragraph{Mechanism: bimodality suppression vs.\ peak precision.}
The mechanism behind the variance-vs-peak-precision tradeoff
is visible in three interlocking signatures.  First, the
standard deviation ratio of $3.21\times$ exceeds the
$\sqrt{3}\approx 1.73\times$ reduction expected from pure
variance reduction over $N=3$ independent samples; the excess
factor of $1.86\times$ indicates that median aggregation does
not merely smooth Gaussian noise but actively suppresses
single-sample bimodality.  Second, the mean--median
divergence in arm~B (mean $27.5$, median $8.4$) versus the
mean--median agreement in arm~A (mean $18.3$, median
$19.95$) is the canonical signature of bimodal vs.\ unimodal
distribution.  Third, the threshold-stratified hit rate
inverts at the $\leq\!5$~px boundary: arm~B captures $50\%$
of cases under sub-pixel precision (single-sample lucky),
while arm~A captures only $10\%$ but achieves $100\%$ under
the relaxed $\leq\!50$~px threshold (outlier-suppression
benefit).

\paragraph{Practical implication: threshold-dependent aggregation.}
The choice of aggregation strategy is therefore
target-threshold-dependent rather than universally better or
worse.  Tasks demanding sub-$5$~px precision benefit from
$N=1$'s winner-take-all property, accepting the $30$\,pp
risk of gross-error outliers ($>\!50$~px tail) for the
$40$\,pp gain in sub-pixel precision (single-sample lucky
alignment).  Tasks requiring stable $\leq\!50$~px
localization benefit from $N=3$ median's outlier
suppression, sacrificing peak precision but eliminating
catastrophic failures.  This positions aggregation strategy
as a tunable design parameter rather than a fixed pipeline
choice; the variance-vs-peak-precision frontier traces a
Pareto curve that an application can place itself on by
selecting~$N$.

\paragraph{Methodological positioning: orthogonal axis of contribution.}
This finding extends the paper's contribution along an axis
orthogonal to the deterministic-routing claim of
Section~\ref{sec:discussion-stage-c}: while routing is the
\emph{geometric} backbone of the pipeline, aggregation is
its \emph{statistical} backbone.  Together they constitute
a co-design space whose joint tradeoffs depend on
application-level accuracy thresholds.  We note that the
variance-vs-peak-precision tradeoff is model-agnostic: the
aggregation analysis applies to any phase-1 sample
distribution and is not specific to the Opus--Sonnet
contrast that originally motivated the W3 protocol
(Section~\ref{sec:results}).  The result thereby refines an
earlier ``model-specific drift'' framing into a
methodological contribution that travels with the pipeline
rather than with the model.

\paragraph{Third regime: arm~C N=3 plurality-vote on tool-use routing.}
The arm~C W3 Opus protocol ($N\!=\!3$ plurality-vote
aggregation on the LLM-routed \texttt{declare\_target} CLI
path, 30 cases completed across three stages 2026-04-28/29)
provides a third regime intersecting the routing-path and
aggregation-strategy axes.  The aggregate $\leq\!20$\,px hit
rate of $\mathbf{86.7\%}$ on the pinpoint method exceeds the
arm~A baseline ($73.3\%$, heuristic rule with $N\!=\!3$
median) by $+13.4$\,pp, completing a reversal of the
Stage~1 mid-fire deficit picture and providing the
strongest single-arm precision result in the W3 series.

\paragraph{Mode A/B partition reveals the selectivity gate.}
Within the \texttt{precision} category (10 sub-pixel target
cases), arm~C performance partitions into two modes
\emph{by target geometry}, not by sub-pixel size.
\textbf{Mode~A (point-target, 6 cases)}---red pixel, blue
needle tip, green dot, purple diamond, yellow line endpoint,
arrow nose---produces $\mathbf{6/6\!=\!100\%}$ hit at
$\leq\!20$\,px with mean $4.13$\,px.  \textbf{Mode~B
(boundary-ambiguous, 4 cases)}---orange-circle center,
crosshair intersection, white-square top-left corner,
innermost-ring center---produces $3/4\!=\!75\%$ hit, with the
crosshair-intersection ($30.8$\,px) as the sole pinpoint miss.
The two-mode partition isolates the arm~C selectivity gate:
$N\!=\!3$ plurality-vote on tool-use routing absorbs the
arm~B single-sample bimodality
($\sigma_B\!=\!32.55$ vs.\ $\sigma_A\!=\!10.15$, $3.21\!\times$
ratio exceeding $\sqrt{3}$) when the target geometry admits a
unique convergent vote, and degrades to marginal improvement
when the target's geometric center is itself ambiguous.

\paragraph{Routing path $\times$ aggregation strategy as a non-additive interaction.}
The 2$\times$2 factorial (routing: heuristic versus
LLM-routed; aggregation: $N\!=\!1$ versus $N\!=\!3$) with
arms A (heuristic, $N\!=\!3$), B (LLM-routed, $N\!=\!1$),
and~C (LLM-routed, $N\!=\!3$) leaves the heuristic-$N\!=\!1$
cell unmeasured but is informative as-is: the heuristic
arm~A produces approximately uniform pinpoint mean
($12.1$--$18.3$\,px range across categories) but non-uniform
threshold-stratified hit rates ($50$--$90\%$ range, with
\texttt{precision} as the weakest at $50\%$), whereas the
LLM-routed arms B,~C exhibit category- and
geometry-dependent variance (Mode~A/B partition above).  The benefit of
$N\!=\!3$ aggregation is therefore \emph{conditioned} on
the routing path: arm~C's $+13.4$\,pp advantage over arm~A
is realised primarily on Mode~A point-targets where
LLM routing both has signal to extract and converges on
a vote, and the gain is substantively reduced on Mode~B
boundary-ambiguous targets.  Notably, the \texttt{precision}
category---where arm~A is also at $50\%$ hit, matching
arm~B---shows the strongest cross-regime jump to arm~C's
$90\%$, which is suggestive of non-additive interaction
between routing path and aggregation strategy; however, with
the (heuristic, $N\!=\!1$) cell unmeasured, a clean factorial
interaction-effect estimate cannot be made from this data.
The pattern is consistent with $N\!=\!3$ plurality-vote on
tool-use routing occupying a regime not directly accessible
to $N\!=\!3$ heuristic or $N\!=\!1$ tool-use under the
partial factorial measured here.  This suggestive
non-additivity indicates that
aggregation-strategy and routing-path choices
should not be tuned in isolation; the selectivity gate
identified here is a candidate target for future
work on adaptive-sampling control.

\subsection{Resolution-method interaction}
\label{sec:resolution-method}

The guardrail set we propose was largely shaped by Claude
Sonnet~4's effective $\sim\!1024$-pixel vision pipeline.
With higher-resolution VLMs (e.g., Claude Opus~4.7's
$\sim\!2576$-pixel input), several future-work
guardrails---G-3 (hybrid post-HIT verification), G-5 (letter
labels), G-6 (hollow ring dots)---may transition from
speculative to empirically validatable: their failure modes
in the guardrail-applied campaign were partially attributable to
resolution-induced ambiguity (close-spaced markers
indistinguishable, crowded toolbar text unreadable) rather
than fundamental methodological flaw.  Conversely, our
$1024$-pixel zoom cap (G-2) is itself resolution-dependent
and should scale with VLM input resolution.  As VLMs continue
to advance---and as displays push beyond 4K toward 8K---the
guardrail--resolution co-design space remains under-explored.
Future work should systematically map guardrail effectiveness
across (model-resolution $\times$ screen-resolution
$\times$ target-size) regimes.
